% ============================================================================
% WORKBOOK — Section 5.6: Graphing Solutions to Inequalities on a Number Line
% CCSS 7.EE.B.4.B
% Practice-focused reinforcement for the study guide topic
% ============================================================================
\section{Graphing Inequalities on a Number Line}
\topicTitle{Graphing Solutions to Inequalities on a Number Line}

\begin{quickReview}{Graphing Inequality Solutions}
To graph the solution of an inequality on a number line:
\begin{itemize}[leftmargin=*]
    \item \textbf{Open circle} ($\circ$): the boundary is \textbf{NOT} included ($<$ or $>$).
    \item \textbf{Closed circle} ($\bullet$): the boundary \textbf{IS} included ($\leq$ or $\geq$).
    \item \textbf{Shade} in the direction of all solutions.
\end{itemize}

\medskip
\textbf{Example:} $x \geq -2$ → Closed circle at $-2$, shade to the right.

$x < 5$ → Open circle at $5$, shade to the left.

\smallskip
\textbf{Memory trick:} $\leq$ and $\geq$ have a line underneath (``equal'') → closed (filled) circle. $<$ and $>$ have no line → open (empty) circle.
\end{quickReview}

% ── Warm-Up ──────────────────────────────────────────────────────────────────
\begin{practiceBox}[funGreen]{\faSun~Warm-Up}

\practiceHeader[funGreen]{Open or Closed? Left or Right?}
For each inequality, state (a) open or closed circle and (b) shade left or right.

\prob $x > 4$ \answerBlank[4cm]
\answerExplain{Open circle at $4$, shade right}{The symbol $>$ does not include $4$, so use an open circle. Solutions are greater than $4$, so shade to the right.}

\prob $n \leq -1$ \answerBlank[4cm]
\answerExplain{Closed circle at $-1$, shade left}{The symbol $\leq$ includes $-1$, so use a closed circle. Solutions are $-1$ or less, so shade to the left.}

\prob $y \geq 0$ \answerBlank[4cm]
\answerExplain{Closed circle at $0$, shade right}{The symbol $\geq$ includes $0$, so use a closed circle. Solutions are $0$ or greater, so shade to the right.}

\prob $t < 7$ \answerBlank[4cm]
\answerExplain{Open circle at $7$, shade left}{The symbol $<$ does not include $7$, so use an open circle. Solutions are less than $7$, so shade to the left.}

\prob $a \geq -3$ \answerBlank[4cm]
\answerExplain{Closed circle at $-3$, shade right}{The symbol $\geq$ includes $-3$, so use a closed circle. Solutions are $-3$ or greater, so shade to the right.}

\end{practiceBox}

% ── Practice A: Solve and Describe the Graph ─────────────────────────────────
\begin{practiceBox}{Solve and Describe the Graph}

\practiceHeader{Solve the inequality, then describe how you would graph it.}

\prob $x + 6 > 10$ \answerBlank[4cm]
\answerExplain{$x > 4$; open circle at $4$, shade right}{Subtract $6$: $x > 4$. Since $x$ is strictly greater than $4$, use an open circle at $4$ and shade to the right.}

\prob $3n - 9 \leq 6$ \answerBlank[4cm]
\answerExplain{$n \leq 5$; closed circle at $5$, shade left}{Add $9$: $3n \leq 15$. Divide by $3$: $n \leq 5$. Since $5$ is included, use a closed circle at $5$ and shade to the left.}

\prob $-2y > 14$ \answerBlank[4cm]
\answerExplain{$y < -7$; open circle at $-7$, shade left}{Divide by $-2$ and flip the sign: $y < -7$. Open circle at $-7$, shade to the left.}

\prob $\dfrac{k}{4} + 2 \geq 5$ \answerBlank[4cm]
\answerExplain{$k \geq 12$; closed circle at $12$, shade right}{Subtract $2$: $\frac{k}{4} \geq 3$. Multiply by $4$: $k \geq 12$. Closed circle at $12$, shade right.}

\prob $5 - x < 8$ \answerBlank[4cm]
\answerExplain{$x > -3$; open circle at $-3$, shade right}{Subtract $5$: $-x < 3$. Multiply by $-1$ and flip: $x > -3$. Open circle at $-3$, shade right.}

\prob $-3a + 1 \geq 10$ \answerBlank[4cm]
\answerExplain{$a \leq -3$; closed circle at $-3$, shade left}{Subtract $1$: $-3a \geq 9$. Divide by $-3$ and flip: $a \leq -3$. Closed circle at $-3$, shade left.}

\end{practiceBox}

% ── Practice B: True or False ────────────────────────────────────────────────
\begin{practiceBox}[funPurple]{True or False?}

\trueOrFalse{The graph of $x < 3$ uses a closed circle at $3$.}
\answerExplain{False}{The symbol $<$ means ``less than'' but not equal to $3$. Since $3$ is not included, we use an open circle, not a closed one.}

\trueOrFalse{The graph of $x \geq -5$ includes the point $-5$.}
\answerExplain{True}{The symbol $\geq$ means ``greater than or equal to.'' Since $-5$ satisfies the ``equal to'' part, it is included (closed circle).}

\trueOrFalse{If $n > 2$, then $n = 2$ is a solution.}
\answerExplain{False}{The symbol $>$ means strictly greater than. The number $2$ is not greater than $2$, so it is not a solution.}

\trueOrFalse{The graph of $y \leq 4$ shades to the right of $4$.}
\answerExplain{False}{The symbol $\leq$ means $y$ is $4$ or less. ``Less'' means to the left on the number line, so the graph shades to the left.}

\end{practiceBox}

% ── Practice C: Is It a Solution? ────────────────────────────────────────────
\begin{practiceBox}[funOrange]{Is the Number a Solution?}

\practiceHeader[funOrange]{Check whether the given value is a solution. Write Yes or No.}

\begin{multicols}{2}
\prob Is $x = 3$ a solution to $x > 3$? \answerBlank[2cm]
\answerExplain{No}{$3 > 3$ is false. The number $3$ is not strictly greater than $3$.}

\prob Is $x = 5$ a solution to $x \geq 5$? \answerBlank[2cm]
\answerExplain{Yes}{$5 \geq 5$ is true because $5$ is equal to $5$ (the ``or equal to'' part).}

\prob Is $x = -4$ a solution to $2x + 1 < -5$? \answerBlank[2cm]
\answerExplain{Yes}{Substitute: $2(-4) + 1 = -8 + 1 = -7$. Is $-7 < -5$? Yes, so $x = -4$ is a solution.}

\prob Is $x = 0$ a solution to $-3x \geq 6$? \answerBlank[2cm]
\answerExplain{No}{Substitute: $-3(0) = 0$. Is $0 \geq 6$? No, so $x = 0$ is not a solution.}
\end{multicols}
\end{practiceBox}

% ── Practice D: Word Problems ────────────────────────────────────────────────
\begin{practiceBox}[funBlue]{\faGlobe~Word Problems}

\wordProblem{A school bus holds at most $48$ students. There are already $31$ students on the bus. Write and solve an inequality for how many more students $s$ can board. Describe the graph.}{students}
\answerExplain{$s \leq 17$; closed circle at $17$, shade left toward $0$}{Inequality: $s + 31 \leq 48$. Subtract $31$: $s \leq 17$. At most $17$ more students can board. Graph: closed circle at $17$, shade left (but $s \geq 0$ in context).}

\wordProblem{A cell phone plan charges $\$0.10$ per text plus $\$15$ per month. You want to spend less than $\$25$. How many texts can you send?}{texts}
\answerExplain{Fewer than $100$ texts ($t < 100$); open circle at $100$, shade left}{Inequality: $0.10t + 15 < 25$. Subtract $15$: $0.10t < 10$. Divide by $0.10$: $t < 100$. Open circle at $100$, shade left.}

\wordProblem{A carpenter needs to cut boards that are at least $3$ feet long from a $10$-foot board. Each cut wastes $0.5$ feet. He has already made $2$ cuts. Write and solve an inequality for how many more cuts $c$ he can make, and describe the graph.}{cuts}
\answerExplain{$c \leq 12$; closed circle at $12$, shade left}{Remaining board after $2$ cuts: $10 - 2(0.5) = 9$ ft. Each additional cut uses $0.5$ ft. He needs at least $3$ ft: $9 - 0.5c \geq 3$. Subtract $9$: $-0.5c \geq -6$. Divide by $-0.5$ and flip: $c \leq 12$.}

\end{practiceBox}

% ── Challenge ────────────────────────────────────────────────────────────────
\begin{challengeBox}

\prob Write an inequality whose graph has a closed circle at $-2$ and shading to the right. Then give three values that are solutions and one that is not.
\answerExplain{$x \geq -2$; solutions: $-2$, $0$, $5$; non-solution: $-3$}{A closed circle at $-2$ shading right means $x \geq -2$. Any number $-2$ or greater is a solution. For example, $-2$, $0$, and $5$ all satisfy $x \geq -2$. The number $-3$ is less than $-2$, so it is not a solution.}

\prob A student claims: ``If $x < -1$, then $x^{2} > 1$.'' Is this always true? Explain and give an example.
\answerExplain{True}{If $x < -1$, then $|x| > 1$. Squaring a number with absolute value greater than $1$ always gives a result greater than $1$. For example: $x = -3$, $(-3)^2 = 9 > 1$. Or $x = -1.5$, $(-1.5)^2 = 2.25 > 1$.}

\end{challengeBox}

\encouragement{Wonderful job graphing inequalities! You can now solve and visualize solutions like a pro!}
